% sections/conclusion.tex

We have presented a stochastic model of haematopoietic lineage commitment that
reconciles two bodies of experimental evidence that have appeared incompatible
for three decades. At single-cell molecular copy numbers, erythroid commitment
is a stochastic event driven by noise in the transcription decision layer, not
by cytokine dose. At population copy numbers, the same circuit produces the
sharp, cytokine-dependent dose-response that generated the instructive tradition.
Both bodies of evidence are simultaneously correct because they measure the same
circuit at different scales.

The GCSF/EPO ratio sweep extends this finding to the two-cytokine control
surface: the true bifurcation threshold is controlled by receptor occupancy
ratio, not by the absolute dose of either signal, and it is abrupt. EPO
dose-response experiments that find flat commitment fraction are measuring the
system with an uninformative measurement window, not probing an EPO-insensitive
circuit. The three-layer cascade architecture explains why all three layers are
EPO-dose-insensitive by distinct mechanisms, and why EPO therefore functions
as a binary permissive signal rather than a probability rheostat.

The cytokine-withdrawal and pulse-duration experiments define the boundary
between cytokine-dependent priming and irreversible commitment: the model has
the former but not the latter, and the absence of an epigenetic layer is the
clearest structural gap. The non-monotone pulse-duration result---stable
post-withdrawal erythroid fate achievable only within a narrow, bracketed window,
with both shorter and longer exposures reverting by distinct mechanisms---is the
single most experimentally non-obvious prediction from this work.

The largest dataset reported here, the factorial EPO--pH sweep across 2400
stochastic replicates, establishes nuclear pH as a quantitatively significant
modulator of the EPO commitment threshold in a direction and by a magnitude that
mean-field analysis cannot predict. More than half of the conditions showed
genuine stochastic bistability, and the pH-dependent shift in EPO* carries direct
implications for the interpretation of EPO hyporesponsiveness in metabolic
acidosis. A single EPO concentration near the bifurcation produces opposite
majority fates at the two physiological pH extremes: this is an experimentally
testable prediction with a binary, easily scored outcome.

The Signal Hierarchical Petri Net formalism that enabled these simulations is
described in the final section of this paper as a reusable resource. Five
experimentally falsifiable predictions are available to guide future single-cell
wet-lab studies; together they constitute a roadmap for validating or refuting
the mechanistic claims made here at single-cell resolution.
